\section{Run Ledger}
\label{app:ledger}

The main matrix is $8 \times 4 \times 5 \times 3 = 480$ runs (8 datasets, 4 model families, lookbacks $L \in \{96, 336, 720, 1440, 2880\}$, seeds $\{2021, 2022, 2023\}$). \textbf{All 480 runs completed}, and every cell in the paper's tables aggregates three seeds. The saturation probe (electricity and traffic, $L \in \{2880, 5760\}$, 2 model families, 3 seeds) completed all 24 of its verdict-bearing runs, plus 16 single-seed reference cells at $L \le 1440$.

\subsection{Provenance: the initially missing 135 runs}

The matrix was not complete at the first snapshot. Of the 480 planned runs, 135 were missing at that point, for execution-environment reasons only; every one was subsequently re-run to completion on idle capacity, and no completed run was ever excluded from the analysis. We keep the original failure classification here for provenance (Table \ref{tab:ledger}):

\begin{itemize}
\item \textbf{GPU memory failures under co-placement (46 runs).} 38 runs raised \texttt{OutOfMemoryError} and 8 raised \texttt{torch.AcceleratorError}, which inspection of the logs shows to be the same CUDA out-of-memory condition reported asynchronously. These failures concentrated exactly where memory pressure was: all 46 were on the three wide datasets (traffic 23, electricity 20, weather 3), and 35 of 46 were PatchTST or TimesNet runs, the two families whose memory footprint grows with channel count and window length. They spanned all lookbacks ($L{=}96$: 9, 336: 10, 720: 8, 1440: 9, 2880: 10), i.e.\ they reflected co-placement contention rather than a lookback effect. Every one of these runs completed on its first attempt once re-run with one job per GPU.
\item \textbf{Time-box cut, never started (65 runs).} When the scheduling time box closed, the remaining queued runs were cut before launch; all 65 were third-seed ($2023$) runs, spread evenly across the four model families (16--17 each).
\item \textbf{Time-box cut, killed mid-run (24 runs).} Arms in flight at the cut (8 with seed 2021, 7 with 2022, 9 with 2023).
\end{itemize}

\begin{table}[t]
\caption{Classification of the 135 initially missing main-matrix runs (all since re-run to completion). The memory class merges \texttt{OutOfMemoryError} (38) and asynchronously reported CUDA OOM surfacing as \texttt{AcceleratorError} (8). All launch configurations are included in the released repository.}
\label{tab:ledger}
\centering\small
\begin{tabular}{llr}
\toprule
Cause class & Where it concentrated & Arms \\
\midrule
GPU memory (OOM $+$ async CUDA OOM) & wide datasets: traffic 23, electricity 20, weather 3 & 46 \\
\quad of which PatchTST / TimesNet & wide-data $\times$ heavy-model cells & 35 \\
\quad of which iTransformer / DLinear & wide datasets only & 11 \\
Time-box cut, never started & all seed 2023 (third-seed runs) & 65 \\
Time-box cut, killed mid-run & mixed seeds (8/7/9 for 2021/2022/2023) & 24 \\
\midrule
Total & & 135 \\
\bottomrule
\end{tabular}
\end{table}

\subsection{Consequences for the reported conclusions}

Completing the matrix moved two cell argmins, both in DLinear and both on cells that previously carried two seeds: ETTh2 ($1440 \to 720$) and ETTm2 ($1440 \to 2880$). Neither touches the paper's ordering claims (the short/long poles are unchanged on every family) or the two-axis analysis; both cells' new values are reflected in Table \ref{tab:bestl}, the predictability targets of Section \ref{sec:predictfull}, and every downstream number. The zero-shot sweep has 8 open runs at $L{=}2880$ on the six non-wide datasets, which are moot because Chronos-Bolt truncates contexts to 2048 steps.
